\scanpage{146}
\mspage{4}
sur $C'_{\unsure{0|o}}$, \deleted{\unsure{est}} donne un morphisme d'adjonction
\[
 \deleted{\unsure{A}}\,\underset{\varprojlim_e F}{\underset{\Vert}{F}}\longrightarrow\alpha_*(\alpha^* F)=\varprojlim_{C'}\alpha^*F
\]
\deleted{\unsure{que}} qui est un iso. Et l'axiome $\mathrm{Dér}_{5}$ \deleted{affirme} \inserted{\unsure{exprime|assure}} qu'il doit en être ainsi si $C'$ est contractile. (N.B.
\margin{?}
\deleted{Cela élimine les dérivateurs de nature triviale, définis \deleted{\unsure{par les}} par des cat. de modèles et ``pointés'', sans localisation par un ens. de flèches $\Sigma$. Par exemple la catégorie $C':$ (}
\diagram{\begin{tikzcd}[ampersand replacement=\&,row sep=tiny,column sep=tiny]
\&\deleted{\cdot}\arrow[d]\\\deleted{\cdot}\arrow[r]\&\deleted{\cdot}
\end{tikzcd}}
\deleted{) (\unsure{avec un|avec|un} objet final mais pas initial, donc pas \deleted{\unsure{cofil}} filtr. à g.) \deleted{\unsure{doit nous donner une}} donne un foncteur $C'\unsure{-|\to}e$ ``anodin'' cohomologiquement, mais pas anodin pour $\varprojlim$)}

la condition \deleted{de cofinalité} filtrante dans la proposition \deleted{\unsure{est}} est \inserted{d'ailleurs} une condition \uline{suffisante} d'anodinité mais nullement néc. Si p. ex. $C$ est la catégorie finale, alors la condition d'anodinité pour $C'-e$ est simplement que $C'$ soit connexe non vide, \deleted{\unsure{ou}} i.e. $\pi_0(C')=\{0\}$. (Condition qui visiblement néglige les $\pi_i$ supérieurs)

Soit $\boxed{\alpha:X'-X}$ un foncteur de petites catégories. La condition suggérée par la prop., que les $\,{}_{x_0\backslash}X'$ soient contractiles, signifie que $\alpha$ est une hot-équivalence co-``stricte, i.e. que $\alpha^\circ:X'{}^\circ\to X^\circ$ est une hot-équiv ``localement sur $X^\circ$'', donc aussi globalement. Si \deleted{\unsure{$\alpha$}}
% archive sheet number: 146
